\documentclass[border=7bp]{standalone}
% Shared preamble for the report's editable vector figures.
\usepackage{fontspec}
\usepackage{xeCJK}
\setmainfont{Calibri}
\setCJKmainfont{SimSun}[AutoFakeBold=1.5]
\usepackage{amsmath}
\usepackage{tikz}
\usetikzlibrary{arrows.meta,positioning,calc,shapes.geometric,fit}
\definecolor{Ink}{HTML}{25313D}
\definecolor{Blue}{HTML}{315E80}
\definecolor{PaleBlue}{HTML}{E7EFF5}
\definecolor{Muted}{HTML}{626B74}
\definecolor{PaleGray}{HTML}{F0F1F3}
\tikzset{
  >=Stealth,
  every node/.style={font=\fontsize{10bp}{13bp}\selectfont,text=Ink},
  flow/.style={->,draw=Ink,line width=.65pt},
  frame/.style={rectangle,draw=Ink,line width=.65pt,align=center,
    minimum height=.85cm,text width=5cm,inner sep=5pt},
  emphasis/.style={frame,draw=Blue,fill=PaleBlue},
  supporting/.style={frame,draw=Muted,fill=PaleGray},
  annotation/.style={font=\fontsize{9bp}{12bp}\selectfont,text=Muted,align=center},
}

\begin{document}
\begin{tikzpicture}
\path (-7.3,0) rectangle (7.3,0);

\node[emphasis,text width=11cm] (goal) at (0,0) {输入目标棋盘与不相交将牌分组};
\node[frame,text width=11cm] (abstract) at (0,-1.4) {构造抽象状态：保留本组将牌和空位的位置\\移动本组牌收费 1，移动其他牌收费 0};
\node[frame,text width=4.2cm] (cold) at (-2.9,-3.15) {冷缓存：从抽象目标反向建表\\0-1 BFS 求抽象最短距离};
\node[frame,text width=4.2cm] (warm) at (2.9,-3.15) {热缓存：核对目标、分组与文件\\匹配后加载距离表};
\node[supporting,text width=11cm] (table) at (0,-5) {准备完成：距离以紧凑数组保存，抽象状态排名作为下标};
\node[emphasis,text width=11cm] (query) at (0,-6.5) {搜索时：棋盘 → 每组抽象编码 → 查表 → 距离求和\\$h_{\mathrm{PDB}}(s)=\sum_i d_i(s)$};
\node[frame,text width=11cm] (use) at (0,-8) {A* 按 $g+h$ 选取候选；IDA* 用 $g+h$ 与阈值比较};
\draw[flow](goal)--(abstract);\draw[flow](abstract.south)-|(cold.north);\draw[flow](abstract.south)-|(warm.north);
\coordinate (join) at (0,-4.3);
\draw[draw=Ink,line width=.65pt](cold.south)|-(join);
\draw[draw=Ink,line width=.65pt](warm.south)|-(join);\draw[flow](join)--(table.north);
\draw[flow](table)--(query);\draw[flow](query)--(use);
\node[annotation,text width=11cm] at (0,-9) {建表与加载属于准备成本，查询属于搜索成本；可加性来自不重复的动作成本划分};

\end{tikzpicture}
\end{document}
